\chapter{1954:Linus Pauling}

\label{chap:chemistry1954}

The Nobel Prize in Chemistry for the year 1954 was awarded to Linus Pauling.

\textbf{Motivation:}

for his research into the nature of the chemical bond and its application to the elucidation of the structure of complex substances

\section{Linus Pauling}

\subsection*{Early Life and Education}

Linus Pauling was born on 1901-02-28 in Portland, Oregon, USA.

\subsection*{Career and Major Research}

Pauling studied chemical engineering at Oregon Agricultural College and took his doctorate at the California Institute of Technology in 1925. He spent most of his career at Caltech, where he became professor of chemistry. He applied quantum mechanics to chemistry, introducing the concepts of hybrid orbitals, electronegativity, and resonance, which he presented systematically in \emph{The Nature of the Chemical Bond} (1939). He also used X-ray crystallography to determine the structures of minerals and, later, the structure of the alpha helix in proteins.

\subsection*{Nobel Prize-winning Work}

The work for which Linus Pauling was awarded the Nobel Prize in Chemistry in 1954 was described by the Nobel Committee as: "for his research into the nature of the chemical bond and its application to the elucidation of the structure of complex substances".

Pauling showed that the properties of chemical bonds, including bond lengths, bond angles, and energies, could be understood and predicted from quantum mechanical principles. His concept of resonance described molecules such as benzene as hybrids of contributing structures, and his scale of electronegativity explained the polarity of bonds between different elements.

\subsection*{Significance and Impact}

The Nature of the Chemical Bond became the standard account of chemical bonding and shaped the teaching of chemistry for decades. Pauling's methods were used to determine the structures of complex substances from minerals to proteins, and his alpha helix proposal was fundamental to the later solution of protein structures.

\subsection*{Later Career and Honors}

Pauling left Caltech in the 1960s and later worked at the Center for the Study of Democratic Institutions in Santa Barbara and Stanford University. In 1962 he was awarded the Nobel Peace Prize for his campaign against nuclear weapons testing, making him the only person to receive two unshared Nobel Prizes. He died on 1994-08-19 in Big Sur, California.

\subsection*{Legacy}

Pauling is remembered as one of the most influential chemists of the twentieth century, whose work defined modern structural chemistry. His name survives in concepts taught to every student of chemistry, and his double Nobel Prize remains unique.

\subsection*{Modern Developments}

Pauling's work on the nature of the chemical bond, including his concept of hybrid orbitals and electronegativity, remains the foundation of modern quantum chemistry and the understanding of molecular structure. His later identification of the alpha-helix and beta-sheet structures of proteins founded structural biology and protein science. Pauling's legacy also extends to medicine, where his ideas shaped research on vitamins and disease, and to the broader application of chemistry to biology.

\subsection*{Selected Publications}

\begin{itemize}

\item \emph{The Nature of the Chemical Bond}, Cornell University Press, 1939.

\end{itemize}

\clearpage

\section*{Chapter Summary}

The 1954 prize recognized research into the nature of the chemical bond. Linus Pauling explained bonding with quantum mechanics, introduced concepts such as electronegativity and resonance, and applied them to the structures of complex substances, work summed up in \emph{The Nature of the Chemical Bond} and later honored with the 1962 Nobel Peace Prize.

